\documentclass[11pt,a4paper]{article}

% Pakete
\usepackage[utf8]{inputenc}
\usepackage[T1]{fontenc}
\usepackage[ngerman]{babel}
\usepackage[left=15mm,right=15mm,top=15mm,bottom=15mm]{geometry}
\usepackage{helvet}
\renewcommand{\familydefault}{\sfdefault}
\usepackage{amsmath,amssymb}
\usepackage{tikz}
\usetikzlibrary{arrows.meta,positioning,patterns,decorations.pathmorphing}
\usepackage{tcolorbox}
\tcbuselibrary{skins}
\usepackage{multicol}

% Farben (sparsam, DDR-Stil)
\definecolor{rahmen}{HTML}{000000}
\definecolor{hintergrund}{HTML}{f5f5f5}

% Keine Seitenzahl
\pagestyle{empty}

% Keine Einrückung
\setlength{\parindent}{0pt}
\setlength{\parskip}{4pt}

\begin{document}

% ════════════════════════════════════════════════════════════════
% TITEL
% ════════════════════════════════════════════════════════════════
\begin{center}
\fbox{\parbox{0.95\textwidth}{\centering\Large\textbf{DOPPELSPALT MIT ELEKTRONEN}\\[2pt]\normalsize Welle-Teilchen-Dualismus | Komplementarität}}
\end{center}

\vspace{3mm}

% ════════════════════════════════════════════════════════════════
% ZWEI SPALTEN
% ════════════════════════════════════════════════════════════════
\begin{multicols}{2}

% ────────────────────────────────────────────────────────────────
% VERSUCHSAUFBAU
% ────────────────────────────────────────────────────────────────
\textbf{1. VERSUCHSAUFBAU}

\begin{center}
\begin{tikzpicture}[scale=0.65, every node/.style={font=\footnotesize}]
    % Quelle
    \draw[thick, fill=gray!20] (0,0) circle (0.5);
    \node at (0,0) {e$^-$};
    \node[below] at (0,-0.7) {Quelle};

    % Pfeil
    \draw[->, thick] (0.7,0) -- (2.3,0);

    % Doppelspalt
    \fill[gray!70] (2.8,-2) rectangle (3.1,2);
    \fill[white] (2.8,-0.7) rectangle (3.1,-0.3);
    \fill[white] (2.8,0.3) rectangle (3.1,0.7);
    \node[below] at (2.95,-2.2) {Doppelspalt};

    % Schirm
    \fill[gray!30] (6.5,-2) rectangle (6.8,2);
    \node[below] at (6.65,-2.2) {Schirm};

    % Interferenzmuster
    \foreach \y in {-1.5,-0.75,0,0.75,1.5} {
        \fill[black!70] (6.85,\y-0.12) rectangle (7.1,\y+0.12);
    }
\end{tikzpicture}
\end{center}

\vspace{2mm}

% ────────────────────────────────────────────────────────────────
% BEOBACHTUNGEN
% ────────────────────────────────────────────────────────────────
\textbf{2. BEOBACHTUNGEN}

\vspace{1mm}
\fbox{\parbox{0.93\linewidth}{
\textbf{Viele Elektronen:} Interferenzmuster\\[2pt]
$\rightarrow$ typisch für \textbf{Wellen}
}}

\vspace{2mm}

\fbox{\parbox{0.93\linewidth}{
\textbf{Einzelne Elektronen:}\\
• Landen als \textbf{Punkte} (Teilchen!)\\
• Aber: Punkte bilden nach und nach\\
\hspace*{6mm}ein \textbf{Interferenzmuster} (Welle!)
}}

\vspace{2mm}

% ────────────────────────────────────────────────────────────────
% DAS PARADOX
% ────────────────────────────────────────────────────────────────
\textbf{3. DAS PARADOX}

\begin{center}
\begin{tikzpicture}[scale=0.6, every node/.style={font=\footnotesize}]
    % Ein Spalt
    \node[align=center] at (0,1.5) {\textbf{1 Spalt offen}};
    \fill[gray!70] (0,0) rectangle (0.3,1);
    \fill[white] (0,0.35) rectangle (0.3,0.65);
    \draw[->, thick] (0.5,0.5) -- (1.5,0.5);
    \node[align=center] at (2.3,0.5) {keine\\Interferenz};

    % Zwei Spalte
    \node[align=center] at (0,-1.5) {\textbf{2 Spalte offen}};
    \fill[gray!70] (0,-2.8) rectangle (0.3,-1.8);
    \fill[white] (0,-2.55) rectangle (0.3,-2.35);
    \fill[white] (0,-2.2) rectangle (0.3,-2);
    \draw[->, thick] (0.5,-2.3) -- (1.5,-2.3);
    \node[align=center] at (2.5,-2.3) {Interferenz-\\muster!};
\end{tikzpicture}
\end{center}

\vspace{1mm}

\fbox{\parbox{0.93\linewidth}{
\centering
\textbf{Ein einzelnes Elektron „weiß"\\von beiden Spalten!}
}}

\columnbreak

% ────────────────────────────────────────────────────────────────
% WELLE-TEILCHEN-DUALISMUS
% ────────────────────────────────────────────────────────────────
\textbf{4. WELLE-TEILCHEN-DUALISMUS}

\vspace{2mm}

\begin{tikzpicture}[every node/.style={font=\small}]
    % Rahmen
    \draw[thick] (0,0) rectangle (5.5,2.8);

    % Überschrift
    \node[anchor=north] at (2.75,2.7) {\textbf{Quantenobjekte zeigen:}};

    % Zwei Seiten
    \draw[dashed] (2.75,0) -- (2.75,2.3);

    % Links: Welle
    \node[anchor=north, align=center] at (1.3,2.2) {\underline{Wellenverhalten}\\[4pt]
    bei Ausbreitung\\
    (Interferenz,\\Beugung)};

    % Rechts: Teilchen
    \node[anchor=north, align=center] at (4.2,2.2) {\underline{Teilchenverhalten}\\[4pt]
    bei Messung\\
    (Punkt auf\\Schirm)};
\end{tikzpicture}

\vspace{3mm}

% ────────────────────────────────────────────────────────────────
% KOMPLEMENTARITÄT
% ────────────────────────────────────────────────────────────────
\textbf{5. KOMPLEMENTARITÄT}

\vspace{2mm}

\fbox{\fbox{\parbox{0.88\linewidth}{
\centering
\textbf{Messung „Welcher Spalt?"}\\[4pt]
$\Downarrow$\\[4pt]
\textbf{Interferenzmuster verschwindet!}
}}}

\vspace{3mm}

\begin{tikzpicture}[every node/.style={font=\small}]
    \draw[thick, double] (0,0) rectangle (5.5,1.8);
    \node[anchor=north, align=center] at (2.75,1.65) {
    \textbf{Komplementaritätsprinzip}\\(Niels Bohr)\\[4pt]
    Weg \textbf{oder} Interferenz\\
    \textbf{– nie beides zugleich!}
    };
\end{tikzpicture}

\vspace{3mm}

% ────────────────────────────────────────────────────────────────
% INTERPRETATION
% ────────────────────────────────────────────────────────────────
\textbf{6. INTERPRETATION}

\vspace{1mm}

\fbox{\parbox{0.93\linewidth}{
Das Elektron breitet sich als\\
\textbf{Wahrscheinlichkeitswelle} aus.\\[3pt]
Erst bei der Messung „entscheidet"\\
es sich für einen Ort.
}}

\end{multicols}

\vspace{2mm}
\hrule
\vspace{2mm}

% ════════════════════════════════════════════════════════════════
% FORMEL-ZEILE (Anschluss)
% ════════════════════════════════════════════════════════════════
\begin{center}
\fbox{\parbox{0.6\textwidth}{\centering
\textbf{Vorwissen:} De-Broglie-Wellenlänge\\[3pt]
$\displaystyle \lambda = \frac{h}{p} = \frac{h}{m \cdot v}$
}}
\end{center}

\end{document}
